% !TeX root = ../tesis.tex
%%%%%%%%%%%%%%%%%%%%%%%%%%%%%%%%%%%%%%%%%%%%%%%%%%%%%%%%%%%%%%%%%%%
% NOTA SOBRE ESTA VERSIÓN (#39, acuerdo 4 de la reunión 2026-10-19)
%%%%%%%%%%%%%%%%%%%%%%%%%%%%%%%%%%%%%%%%%%%%%%%%%%%%%%%%%%%%%%%%%%%
\chapter*{Nota sobre esta versión}
\addcontentsline{toc}{chapter}{Nota sobre esta versión}

Esta versión reúne las correcciones acordadas en las revisiones de octubre y
noviembre de 2026. Para quien haya leído una versión anterior, el
Cuadro~\ref{tab:nota-version} resume qué cambió de nombre, de valor o de
criterio, y dónde se explica cada cambio. Las cifras del documento salen de una
sola corrida de \texttt{VFTModel}, la del 2026-10-04; cuando una figura o un
cuadro usa otra, su pie lo indica.

Algunas convenciones se mantienen en todo el texto. «Modo» designa cada
sistema de transporte, y los modos se dividen en masivos y no masivos según su
capacidad (§\ref{sec:metodo}). El Caso~1 es la red actual, el \textit{statu
quo}; los Casos~2 y~3 son el Anillo MB y el Anillo Metro. Cada sigla se
expande la primera vez que aparece, y los términos en inglés van en cursiva.

\begin{table}[H]
\centering
\caption[Equivalencias entre versiones]{Equivalencias entre la versión
anterior y esta versión.}
\label{tab:nota-version}
\small
\begin{tabularx}{\textwidth}{>{\raggedright\arraybackslash}p{2.3cm}
  >{\raggedright\arraybackslash}X >{\raggedright\arraybackslash}X
  >{\raggedright\arraybackslash}p{2.1cm}}
\toprule
\textbf{Tema} & \textbf{Versión anterior} & \textbf{Esta versión} &
\textbf{Dónde} \\
\midrule
Corrida de \texttt{VFTModel}
  & Varias corridas, de abril a septiembre de 2026
  & Corrida única del 2026-10-04 (\texttt{VFTModel} \texttt{4ef4265}); $\Delta E$
    sigue como corrida exploratoria del 2026-09-02
  & Pies de figuras y cuadros; §\ref{subsec:vulnerabilidad-estructural} \\
Ámbito geográfico
  & Texto y figuras con las 141 demarcaciones del conjunto de datos
  & Texto: \zmvm{} de 76 demarcaciones. Figuras: las 141 del conjunto de
    datos. La Clasificación Garibelt se reporta en ámbito metropolitano
  & §\ref{subsec:cobertura-espacial}, §\ref{subsec:limitaciones-datos} \\
Casos del anillo
  & Caso~2: Metro, 35~km/h, estaciones cada 1,200~m. Caso~3: Metrobús,
    22~km/h, cada 600~m
  & Caso~2: Anillo MB, autobús en vía confinada como el Mexibús, 16.3~km/h que bajan
    a unos 14.2~km/h. Caso~3: Anillo Metro, 36~km/h. Los dos comparten las 94
    estaciones
  & §\ref{subsec:anillo-parametros}, §\ref{subsec:corredor-anillar-limitaciones} \\
Referencia de $T$
  & «Referencia CDMX» de 85~min
  & Tiempo de referencia del modelo; rango de 85 a 180~min como supuesto del
    autor
  & §\ref{subsubsec:garibelt-procedimiento} \\
Escala Garibelt
  & Bandas Crítico, Débil, Aceptable e Idóneo; «banda dominante»
  & Intervalos 0.00--0.25, 0.25--0.50, 0.50--0.75 y 0.75--1.00; «intervalo
    dominante»
  & §\ref{subsubsec:garibelt-escala} \\
Nombre de $DI$
  & Índice de Ruta Directa; símbolo $DF$ en los anexos
  & Factor de Desviación, con símbolo $DI$ en todo el documento
  & §\ref{subsec:ruta-directa} \\
Fuerza Capilar de Macro-Hub
  & $FC_h$ en el Anexo F
  & $FC_H$ en todo el documento
  & §\ref{subsec:fuerza-capilar} \\
\bottomrule
\end{tabularx}
\fuentefigura{Elaboración propia.}
\end{table}
